%% kapa: exam=1 round=34 year=2023 period=2 form=A subject=accounting subject_ko=회계학 seq=27 number=67
%% source: Q-Net 감정평가사 기출문제 (cst00309) artlSeq=5211789 file=raw/1008/5211789/2023년도 제34회 감정평가사 1차시험 2교시.pdf
%% conversion: kordoc 4.24.51 -> tools/gen1.py (auto); review=proofread-sonnet; flags=none
\begin{question}{67}
(주)감평은 20×1년 초 투자부동산(내용연수 10년, 잔존가치 ￦0, 정액법상각)을 ￦200,000에 취득하고 원가모형을 적용하였다. (주)감평은 20×2년부터 동 투자부동산에 대하여 공정가치모형을 적용하기로 하였으며 이러한 회계변경은 정당하다. 20×1년 말, 20×2년 말 동 투자부동산의 공정가치는 각각 ￦190,000, ￦185,000이다. 회계변경효과를 반영하여 20×2년 말 작성하는 비교재무제표(20×1년, 20×2년)에 표시될 금액에 관한 설명으로 옳은 것은?
\begin{choices}
  \item 20×1년도 투자부동산(순액)은 ￦180,000이다.
  \item 20×1년도 투자부동산 감가상각비는 ￦0이다.
  \item 20×1년도 투자부동산평가손익은 ￦0이다.
  \item 20×2년도 투자부동산평가이익은 ￦5,000이다.
  \item 20×2년도 투자부동산(순액)은 ￦190,000이다.
\end{choices}
\end{question}
\answer{2}
